% Original pages: 312--316
% Figures: Figures 2--5 placeholders on original pp. 313--315.
% Uncertainty log: VERIFY comments included for equation (12) index lettering and the trailing word split at the end of original p. 316.

energies, they may, of course, be relevant. In these lectures we
will only consider the simplest theories with a single conserved
spinor $Q_\alpha$ and the basic algebra (11).

The algebra (11) has some dramatic consequences. Since $Q_\alpha$ is
hermitian, $\overline{Q}_\beta = Q_\sigma \gamma^0_{\sigma\beta}$; so (11) can be written
% VERIFY: The source scan is faint in equation (12); the dummy summed index on the second $Q$ / $\gamma^0$ factor should be rechecked against the scan if stricter literal fidelity is required.
\begin{equation}
\{Q_\alpha, Q_\sigma\}\gamma^0_{\sigma\beta} = \gamma^\mu_{\alpha\beta} P_\mu
\tag{12}
\end{equation}
If we multiply by $\gamma^0_{\beta\alpha}$, sum over $\beta$ and $\alpha$, and use the facts $(\gamma^0)^2 = 1$,
$\mathrm{Tr}\,\gamma^0 \gamma^\mu = 4\delta^\mu_0$, we get
\begin{equation}
4P_0 = \sum_\alpha \{Q_\alpha, Q_\alpha\}
\tag{13}
\end{equation}
Here $P_0$ is, of course, the Hamiltonian $H$. For any operator $A$,
$\{A, A\} = 2A^2$. Equation (13) is equivalent to
\begin{equation}
H = \frac{1}{2}\sum_\alpha Q_\alpha^2
\tag{14}
\end{equation}
which is one of the keys to understanding supersymmetry.

It follows immediately from equation (14) that if supersymmetry
is not spontaneously broken --- if the $Q_\alpha$ annihilate the vacuum
$|\Omega\rangle$ --- then the energy of the vacuum is zero. If $Q_\alpha |\Omega\rangle = 0$ then
obviously
\begin{equation}
H |\Omega\rangle = \frac{1}{2}\sum_\alpha Q_\alpha^2 |\Omega\rangle = 0
\tag{15}
\end{equation}
If conversely, supersymmetry is spontaneously broken, that is, if
$Q_\alpha |\Omega\rangle \neq 0$, then
\begin{equation}
\langle \Omega | H | \Omega \rangle = \frac{1}{2}\sum_\alpha \langle \Omega | Q_\alpha^2 | \Omega \rangle = \frac{1}{2}\sum_\alpha \left|Q_\alpha |\Omega\rangle\right|^2 > 0
\tag{16}
\end{equation}
so in this case the vacuum energy is positive.\begingroup\renewcommand{\thefootnote}{\fnsymbol{footnote}}\footnotemark[1]\endgroup

\begingroup
\renewcommand{\thefootnote}{\fnsymbol{footnote}}
\footnotetext[1]{Usually, one is free to add a constant to the Hamiltonian, but here the zero of energy is fixed in a natural way by equation(14). When asking what is the energy of the vacuum, we always have in mind the definition (14) of H.}
\endgroup

Combining these remarks, we see that supersymmetry is spontaneously
broken if and only if the energy of the vacuum is greater
than zero. This is often illustrated by the diagram of figure (2).
In figure (2(a)) a scalar field has a vacuum expectation value, possibly
breaking some internal symmetry. However, supersymmetry is
unbroken because the ground state energy --- the minimum of the potential
--- is zero. In figure (2(b)), the expectation value of the
scalar field is zero, but supersymmetry is spontaneously broken because the
energy of the ground state is greater than zero.

If supersymmetry plays any role in nature it must be spontaneously
broken. This follows immediately from the fact that the $Q_\alpha$
have spin 1/2, so acting on any particle they change its spin by $\pm 1/2$:
\begin{equation}
Q_\alpha |\text{spin } s\rangle \sim |\text{spin } s \pm 1/2\rangle
\tag{17}
\end{equation}
Since $Q_\alpha Q_\alpha$ commutes with the Hamiltonian it does not change the mass
of the particle on which it acts. Since we do not observe in nature
the degeneracies among particles of different spin that would be predicted
by (17), supersymmetry must be spontaneously broken if it is
relevant to nature.

% TODO FIGURE: Original Fig. 2 (panels (a) and (b)) appears here on p. 313.
\begin{center}
\refstepcounter{figure}
\begin{minipage}{0.34\textwidth}
\centering
\fbox{\parbox[c][1.15in][c]{0.24\textwidth}{\centering TODO\\Figure placeholder\\panel (a)\\Original p.~313}}
\par\smallskip
(a)
\end{minipage}\hfill
\begin{minipage}{0.34\textwidth}
\centering
\fbox{\parbox[c][1.15in][c]{0.24\textwidth}{\centering TODO\\Figure placeholder\\panel (b)\\Original p.~313}}
\par\smallskip
(b)
\end{minipage}
\par\smallskip
Figure~\thefigure
\end{center}

A puzzle immediately presents itself. Since broken supersymmetry
means a positive vacuum energy, why does not a cosmological
constant arise as soon as a supersymmetric theory is coupled to
gravity? No good answer is known to this question, and it certainly
is one of the most important questions that must be answered. Actually
the positive vacuum energy of global supersymmetry does not necessarily
become a positive cosmological constant after coupling to
gravity, because the coupling to gravity introduces extra terms in
the scalar potential; these extra terms are not positive definite.
While a cancellation can occur,\textsuperscript{5} leaving zero cosmological constant,
such a cancellation apparently depends on a special choice of parameters,
for which there is no known rationale.

Now, the subject of spontaneously broken supersymmetry has a
special flavor, which arises from the fact that (leaving gravity
aside), supersymmetry is spontaneously broken if and only if the
ground state energy is positive. Suppose that in some theory an approximate
calculation gives an effective potential such as the one
shown in fig. (3). In this approximation, the minimum of the potential
is zero, and supersymmetry is unbroken.

% TODO FIGURE: Original Fig. 3 appears here on p. 314.
\begin{wrapfigure}{l}{0.28\textwidth}
\vspace{-0.5\baselineskip}
\refstepcounter{figure}
\centering
\fbox{\parbox[c][1.2in][c]{0.22\textwidth}{\centering TODO\\Figure placeholder\\Original p.~314}}
\par\smallskip
Figure~\thefigure
\vspace{-0.9\baselineskip}
\end{wrapfigure}

Even if the approximation in question is very accurate, one
cannot be certain, just from this result, that supersymmetry really
is unbroken, because the errors in one's calculation may shift the
potential slightly away from $V = 0$ (figure (4)) (the approximate
calculation and the exact result are indicated in figure (4) by the
solid and dotted line, respectively).

If, instead, an approximate calculation shows that supersymmetry
is spontaneously broken, one can be confident in the result if

% TODO FIGURE: Original Figs. 4 and 5 appear here on p. 315.
\begin{center}
\begin{minipage}{0.42\textwidth}
\centering
\refstepcounter{figure}
\fbox{\parbox[c][1.1in][c]{0.30\textwidth}{\centering TODO\\Figure placeholder\\Original p.~315}}
\par\smallskip
Figure~\thefigure
\end{minipage}\hfill
\begin{minipage}{0.42\textwidth}
\centering
\refstepcounter{figure}
\fbox{\parbox[c][1.1in][c]{0.30\textwidth}{\centering TODO\\Figure placeholder\\Original p.~315}}
\par\smallskip
Figure~\thefigure
\end{minipage}
\end{center}

(figure (5)) the calculated vacuum energy $E$ is much bigger than the
uncertainty $\Delta E$ in the calculation.

On the other hand, for ordinary symmetries, a reliable approximation
(like perturbation theory, in a weakly coupled theory) can
reliably indicate whether the symmetry is spontaneously broken.

I have discussed these matters in much more detail in a recent
paper.\textsuperscript{6} In that paper I also described a simple quantum mechanics
model in which tiny quantum effects shift the energy slightly away
from zero, as in figure (4).

\section{SUPERSYMMETRY AND THE LARGE NUMBERS}

Supersymmetry is a very beautiful idea, but I think it is fair
to say that no one knows what mysteries of nature (if any) it should
explain. To fix ideas, I will make some definite assumptions in the
next few lectures about what problems supersymmetry might solve.

Certainly, the ultimate applications may be in a very different area.
Perhaps supersymmetry is spontaneously broken at very high
energies --- energies that may be as high as the Planck mass $10^{19}$ GeV.
We will assume instead in most of our discussion that supersymmetry
is, more or less, within reach of the new generation of accelerators.
What would we like to explain with Supersymmetry? Of the many
possible answers, I will focus on Dirac's ``problem of the large numbers,''\textsuperscript{7}
which nowadays is often called the ``gauge hierarchy problem.''
As posed by Dirac, the question is why the Planck mass $M_{\rm Pl}$ is so much
larger than the nucleon mass $M_N$: $M_{\rm Pl}/M_N \sim 10^{19}$. In a contemporary
version,\textsuperscript{8} one might ask why the mass scale $M_X$ of presumed grand unification
is so much larger than the mass scale $M_W$ of weak interactions:
$M_X/M_W > 10^{13}$. Starting with Dirac, most physicists have believed
that such large numbers should not be postulated arbitrarily
but must have a definite explanation.

We now know that the ``normal'' masses --- the quark, lepton, and
$W$ and $Z$ masses --- are determined by $SU(2) x U(1)$ breaking. Specifically,
they are all proportional to the expectation value $\langle \phi \rangle$ of the
Higgs boson. So the question is really why $\langle \phi \rangle$ is so tiny compared
to the ``large'' masses --- the mass scale of grand unification, or the
Planck mass.

Looking at the standard Higgs potential
\begin{equation}
V(\phi) = \frac{\lambda}{4}\phi^4 - \frac{m^2}{2}\phi^2
\tag{18}
\end{equation}
we have $\langle \phi \rangle = m/\sqrt{\lambda}$, so the real problem is to explain why $m_\phi \ll M_{\rm Pl}$.

Why would supersymmetry be relevant to this?

As we now understand it, bare mass terms for the quarks and
leptons are forbidden by $SU(2) x U(1)$ gauge invariance. Left-handed
quarks and leptons transform as doublets of $SU(2) x U(1)$, but right-

% VERIFY: Original p. 316 ends mid-word at ``right-''; the continuation is on p. 317 and is intentionally not resolved here.
